\documentclass[10pt,letterpaper,landscape]{article}
\usepackage[margin=.48in]{geometry}
\usepackage[T1]{fontenc}
\usepackage{lmodern,xcolor,listings,hyperref}
\definecolor{ink}{HTML}{17283C}
\definecolor{teal}{HTML}{006B66}
\definecolor{ice}{HTML}{ECF4F5}
\hypersetup{colorlinks,urlcolor=teal}
\pagestyle{empty}
\setlength{\parindent}{0pt}
\setlength{\parskip}{3pt}
\renewcommand{\familydefault}{\sfdefault}
\lstset{basicstyle=\ttfamily\footnotesize,backgroundcolor=\color{ice},
 keywordstyle=\bfseries\color{teal},commentstyle=\color{teal},
 morecomment=[l]{//},morekeywords={module,endmodule,logic,input,output,parameter,
 assign,always_comb,always_ff,begin,end,if,else,case,endcase,default,
 typedef,enum,posedge,initial},columns=fullflexible,keepspaces=true,
 showstringspaces=false,aboveskip=5pt,belowskip=5pt,xleftmargin=4pt}
\newcommand{\topic}[1]{\par\medskip{\large\bfseries\color{teal}#1}\par}
\begin{document}
\color{ink}
{\fontsize{28}{32}\selectfont\bfseries SystemVerilog Cheat Sheet}\hfill
{\small DIGITAL DESIGN \textbar{} RISC-V EDITION}\par
\vspace{3pt}{\color{teal}\hrule height 1.2pt}\vspace{8pt}
\small
\begin{minipage}[t]{.315\textwidth}
\topic{Combinational logic}
\begin{lstlisting}
module mux2 #(parameter int W = 8) (
  input logic [W-1:0] a, b,
  input logic sel,
  output logic [W-1:0] y
);
  assign y = sel ? b : a;
endmodule
\end{lstlisting}
Use \verb|assign| for a continuous expression. Use blocking \verb|=| inside \verb|always_comb|.
\begin{lstlisting}
always_comb begin
  y = '0;  // default on every path
  if (en) y = a;
end
\end{lstlisting}
An unassigned path can infer a latch. Give each RTL signal one driver.

\topic{Bits and operators}
\begin{lstlisting}
8'hA5        // 8-bit hex literal
'0           // fill with zeros
a[7:4]       // select four bits
{hi, lo}     // concatenate
{4{sel}}     // replicate four times
&a           // reduction AND
a & b        // bitwise AND
a && b       // logical AND
\end{lstlisting}
\verb|logic| has four states (0, 1, X, Z). Check truncation and mixed signedness. For \verb|>>>| to sign-fill, the left operand must be signed.

\topic{Named port connections}
\begin{lstlisting}
mux2 #(.W(8)) u_mux (
  .a(in0), .b(in1),
  .sel(choose), .y(out)
);
\end{lstlisting}
\end{minipage}\hfill
\begin{minipage}[t]{.315\textwidth}
\topic{Registers}
\begin{lstlisting}
always_ff @(posedge clk) begin
  if (reset)   q <= '0;
  else if (en) q <= d;
end
\end{lstlisting}
Use nonblocking \verb|<=| for clocked state. With neither reset nor enable asserted, the register holds its value. This reset is synchronous. For an asynchronous active-high reset, use \verb|@(posedge clk or posedge reset)|.

\topic{Finite state machines}
\begin{lstlisting}
typedef enum logic {IDLE, RUN} state_t;
state_t state, next_state;

always_ff @(posedge clk)
  if (reset) state <= IDLE;
  else       state <= next_state;

always_comb begin
  next_state = state;
  case (state)
    IDLE: if (start) next_state = RUN;
    RUN:  if (done)  next_state = IDLE;
    default:         next_state = IDLE;
  endcase
end

assign busy = (state == RUN);
\end{lstlisting}
\textbf{Moore} outputs depend on state. \textbf{Mealy} outputs depend on state and inputs.

\topic{Packed vs.\ unpacked}
\begin{lstlisting}
logic [7:0] mem [0:15];
\end{lstlisting}
The packed dimension gives each word 8 bits. The unpacked dimension gives 16 words. \verb|mem[3][7]| selects bit 7 of word 3.
\end{minipage}\hfill
\begin{minipage}[t]{.315\textwidth}
\topic{Self-checking testbench}
\begin{lstlisting}
module mux2_tb;
  timeunit 1ns;
  timeprecision 1ps;
  logic [7:0] a, b, y;
  logic sel;

  mux2 dut (.a(a), .b(b),
            .sel(sel), .y(y));

  initial begin
    a = 8'h12;
    b = 8'hA5;
    sel = 1'b0;
    #1;
    if (y !== a) $fatal(1, "sel=0");

    sel = 1'b1;
    #1;
    if (y !== b) $fatal(1, "sel=1");
    $display("PASS");
    $finish;
  end
endmodule
\end{lstlisting}
Use \verb|!==| to catch unexpected X or Z bits when the expected value is known. The delays and system tasks shown here are for simulation.

\topic{Debugging checklist}
\textbf{Unexpected X.} Check reset, unconnected inputs, and competing drivers.

\textbf{Wrong width.} Check slices, literal sizes, and signedness before the operation.

\textbf{Wrong cycle.} Drive synchronous stimulus away from the sampling edge. Check outputs after nonblocking updates settle.
\end{minipage}

\vfill
{\color{teal}\hrule height .6pt}\vspace{4pt}
{\footnotesize Based on Harris \& Harris,
\href{https://pages.hmc.edu/harris/ddca/ddcarv.html}{\emph{Digital Design and Computer Architecture: RISC-V Edition}}.
\hfill\href{https://pages.hmc.edu/harris/class/e85/DDCArv_Ch4.pdf}{Chapter 4 lecture slides}}
\end{document}
